% Section 4.3, from lectures/L04/appendix/b_exercises.md.

\section{Exercises}
\label{c4:sec:exercises}\label{c4:app:b}

\begin{aside}
\textbf{How to work these.} The Reflection and Calculation exercises are answered on paper, from
the chapter. The Code exercises are programs for the board: build and flash them the way
\secref{c2:app:c} (Microchip Studio) or \secref{c2:app:d} (Linux) describes, or run them in the
simulator of Appendix~\ref{sim:ch} if you have no board yet.

\textbf{How to check your work.} The Reflection and Calculation exercises are answered in
Appendix~\ref{app:answers}. The Code exercises have no printed solution: the board tells you
whether a program does what the exercise describes. Try each exercise before you read its answer.
\end{aside}

% ------------------------------------------------------------------------------------------
\exerciseset{Concepts}
\label{c4:set:1}

\calcexercise{From reading to voltage}
\label{c4:ex:1}
A potentiometer reads \code{768} from \code{adc_read()}. Given a 5\,V reference, what voltage does
that correspond to? Show the calculation.

\reflexercise{The ADC's own prescaler}
\label{c4:ex:2}
Why does the ADC need its own clock prescaler, separate from the CPU's? What would you expect to
happen to the reading's accuracy if the ADC ran at the full 16\,MHz CPU clock instead?

% ------------------------------------------------------------------------------------------
\exerciseset{Registers}
\label{c4:set:2}

\reflexercise{From A0 to A1}
\label{c4:ex:3}
Which single bit would you flip in \code{ADMUX} to read from \code{A1} instead of \code{A0},
assuming the reference voltage stays the same?

\reflexercise{Waiting for \code{ADIF}}
\label{c4:ex:4}
Explain, in your own words, why the code polls the \code{ADIF} flag in a \code{while} loop after
starting a conversion, rather than reading \code{ADC} immediately after setting \code{ADSC}.

% ------------------------------------------------------------------------------------------
\exerciseset{Program}
\label{c4:set:3}

\codeexercise{A threshold switch}
\label{c4:ex:5}
Extend the PWM-brightness example of \secref{c4:a:4} so that, instead of software PWM, it simply
lights the LED solid whenever the potentiometer reading exceeds \code{512}, and turns it off
otherwise, a basic threshold switch with no PWM. Rewrite \code{run_pwm_cycle()} (rename it
too; it's no longer PWM).

\reflexercise{Removing the wait}
\label{c4:ex:6}
What would you expect to happen if the \code{while (0U == (ADCSRA & (1U << ADIF))) {}} wait were
removed entirely from \code{adc_read()}? Would the program crash outright, or misbehave more
subtly? Explain your reasoning.

\codeexercise{Five dimmable LEDs}
\label{c4:ex:7}
Capstone. Extend the software-PWM approach of \secref{c4:a:4} to dim five LEDs from a single
potentiometer, each individually switchable:
\begin{itemize}
  \item LEDs on D6--D10.
  \item Buttons on D2, D3, D11--D13:
    \begin{itemize}
      \item D2/D3 are \code{INT0}/\code{INT1} from Chapter~\ref{c3:ch}.
      \item D11--D13 share a Pin Change Interrupt, also from Chapter~\ref{c3:ch}.
    \end{itemize}
\end{itemize}
Each button shall toggle the PWM state of an LED:
\begin{itemize}
  \item Button D2 toggles LED D6.
  \item D3 toggles D7.
  \item D11 toggles D8.
  \item D12 toggles D9.
  \item D13 toggles D10.
\end{itemize}
Only LEDs currently ``on'' should respond to the potentiometer's duty cycle; ``off'' LEDs stay dark
no matter where the pot is turned.

Reuse \code{adc_read()} from \secref{c4:a:3} for the potentiometer reading and the duty-cycle math
from \secref{c4:a:4}.
\begin{itemize}
  \item Track which of the five LEDs are enabled using a \code{uint8_t} bitmask with bitwise tests.
  \item Drive the on/off timing for all five within one PWM period instead of just one.
  \item The Pin Change Interrupt runs on any change of D11--D13, so its ISR must work out which
    button was just pressed; use the \code{prev}/\code{now} technique from \secref{c3:a:8}.
\end{itemize}
